\documentclass[11pt,a4paper]{article}
\usepackage[utf8]{inputenc}
\usepackage[spanish]{babel}
\usepackage{amsmath, amssymb}
\usepackage{geometry}
\geometry{top=2.5cm, bottom=2.5cm, left=2.5cm, right=2.5cm}

\begin{document}

\section*{Serie de ejercicios seleccionados para preparar examen del capítulo 10}

\begin{itemize}
    \item[\textbf{1.}] El núcleo de un átomo se puede modelar como constituido por protones y neutrones cercanamente empaquetados. Cada partícula tiene una masa de \(1.67 \times 10^{-27}\,\text{kg}\) y un radio del orden de \(10^{-15}\,\text{m}\).
    \begin{enumerate}
        \item[\textbf{(a)}] Utilice este modelo y los datos proporcionados para estimar la densidad del núcleo de un átomo.
        \item[\textbf{(b)}] Compare su resultado con la densidad de un material, por ejemplo, el hierro. ¿Qué sugieren su resultado y la comparación respecto a la estructura de la materia?
    \end{enumerate}
    \vspace*{7cm}
    
    \item[\textbf{2.}] Para el sótano de una nueva casa, se cava un hoyo en el suelo, con lados verticales que bajan \(2.40\,\text{m}\). Una pared de cimiento de concreto se construye horizontal los \(9.60\,\text{m}\) de ancho de la excavación. Esta pared de cimiento mide \(0.183\,\text{m}\) desde el frente del hoyo del sótano. Durante una tormenta, el drenaje de la calle llena el espacio enfrente de la pared de concreto, pero no el sótano detrás de la pared. El agua no se filtra en el suelo de arcilla. Encuentre la fuerza que ejerce el agua sobre la pared de cimiento. En comparación, el peso del agua está dado por \(2.40\,\text{m} \times 9.60\,\text{m} \times 0.183\,\text{m} \times 1000\,\text{kg/m}^3 \times 9.80\,\text{m/s}^2 = 41.3\,\text{kN}\).
    \vspace*{7cm}
\end{itemize}

\newpage

\begin{itemize}
    \item[\textbf{3.}] Un tanque con un fondo plano de área \(A\) y lados verticales se llena de agua con una profundidad \(h\). La presión es \(P_0\) en la superficie superior.
    \begin{enumerate}
        \item[\textbf{(a)}] ¿Cuál es la presión absoluta en el fondo del tanque?
        \item[\textbf{(b)}] Suponga que un objeto de masa \(M\) y densidad menor a la densidad del agua se coloca en el tanque y flota. No se desborda agua. ¿Cuál es el aumento resultante de presión en el fondo del tanque?
    \end{enumerate}
    \vspace*{7cm}
    
    \item[\textbf{4.}] El 21 de octubre de 2001, Ian Ashpole del Reino Unido logró un récord de altura de \(3.35\,\text{km}\) (\(11000\,\text{pies}\)) impulsado por 600 globos llenos con helio. Cada globo lleno tuvo un radio de aproximadamente \(0.50\,\text{m}\) y una masa estimada de \(0.30\,\text{kg}\).
    \begin{enumerate}
        \item[\textbf{(a)}] Estime la fuerza de flotación total sobre los 600 globos.
        \item[\textbf{(b)}] Estime la fuerza neta hacia arriba sobre los 600 globos.
        \item[\textbf{(c)}] Ashpole utilizó paracaídas para su retorno a la Tierra después de que los globos empezaron a reventarse a elevadas altitudes y disminuyera la fuerza boyante. ¿Por qué se reventaron los globos?
    \end{enumerate}
    \vspace*{7cm}
\end{itemize}

\newpage

\begin{itemize}
    \item[\textbf{5.}] El agua que fluye de una manguera de jardín de diámetro \(2.74\,\text{cm}\) llena una cubeta de \(25\,\text{L}\) en \(1.50\,\text{min}\).
    \begin{enumerate}
        \item[\textbf{(a)}] ¿Cuál es la rapidez del agua que sale del extremo de la manguera?
        \item[\textbf{(b)}] Se acopla una boquilla al extremo de la manguera. Si el diámetro de la boquilla es un tercio del diámetro de la manguera, ¿cuál es la rapidez del agua que sale por la boquilla?
    \end{enumerate}
    \vspace*{7cm}
    
    \item[\textbf{6.}] Un legendario niño holandés salvó a Holanda al poner su dedo en un hoyo de \(1.20\,\text{cm}\) de diámetro en un dique. Si el hoyo estaba \(2.00\,\text{m}\) bajo la superficie del Mar del Norte (densidad \(1030\,\text{kg/m}^3\)),
    \begin{enumerate}
        \item[\textbf{(a)}] ¿cuál fue la fuerza sobre su dedo?
        \item[\textbf{(b)}] Si él hubiera sacado el dedo del hoyo, ¿durante qué intervalo de tiempo, el agua liberada llenaría 1 acre de tierra a una profundidad de 1 pie? Suponga que el hoyo mantuvo constante su tamaño.
    \end{enumerate}
    \vspace*{7cm}
\end{itemize}

\newpage

\begin{itemize}
    \item[\textbf{7.}] En flujo ideal, un líquido de \(850\,\text{kg/m}^3\) de densidad se mueve desde un tubo horizontal de \(1.00\,\text{cm}\) de radio a un segundo tubo horizontal de \(0.500\,\text{cm}\) de radio a la misma elevación que el primer tubo. La presión difiere por \(\Delta P\) entre el líquido en un tubo y el líquido en el segundo tubo.
    \begin{enumerate}
        \item[\textbf{(a)}] Encuentre la rapidez de flujo volumétrico como una función de \(\Delta P\).
        \item[\textbf{(b)}] Evalúe la rapidez de flujo volumétrico para \(\Delta P = 6.00\,\text{kPa}\) y \textbf{(c)} \(\Delta P = 12.0\,\text{kPa}\).
    \end{enumerate}
    \vspace*{7cm}
    
    \item[\textbf{8.}] Una aguja hipodérmica mide \(3.00\,\text{cm}\) de largo y \(0.300\,\text{mm}\) de diámetro. ¿Qué diferencia de presión entre la entrada y la salida de la aguja es necesaria para que el caudal de agua a través de ella sea de \(1.00\,\text{g/s}\)? (Use \(1.00 \times 10^{-3}\,\text{Pa}\cdot\text{s}\) como la viscosidad del agua.)
    \vspace*{7cm}
\end{itemize}

\newpage

\begin{itemize}
    \item[\textbf{9.}] Un avión tiene una masa de \(1.60 \times 10^4\,\text{kg}\) y cada ala tiene un área de \(40.0\,\text{m}^2\). Durante vuelo a nivel, la presión sobre la superficie inferior del ala es \(7.00 \times 10^4\,\text{Pa}\).
    \begin{enumerate}
        \item[\textbf{(a)}] Suponga que el empuje sobre el avión solo se debió a una diferencia de presión. Determine la presión sobre la superficie superior del ala.
        \item[\textbf{(b)}] Más real, una parte significativa del empuje se debe a la deflexión del aire hacia abajo causada por el ala. ¿La presión en (a) es más alta o más baja si se incluye esta fuerza? Explique.
    \end{enumerate}
    \vspace*{7cm}
    
    \item[\textbf{10.}] Hace décadas se creía que los grandes dinosaurios herbívoros, como el \textit{Apatosaurus} y el \textit{Brachiosaurus}, habitualmente caminaban sobre el fondo de los lagos, y sus largos cuellos les permitían tomar aire de la superficie para respirar. El \textit{Brachiosaurus} tenía sus orificios nasales en la parte superior de su cabeza. En 1977, Knut Schmidt-Nielsen puntualizó que para tal criatura sería muy difícil el proceso respiratorio. Como un simple modelo, considere una muestra de \(10.0\,\text{L}\) de aire a presión absoluta de \(2.00\,\text{atm}\), con densidad de \(2.40\,\text{kg/m}^3\), localizada en la superficie de un lago de agua dulce. Encuentre el trabajo requerido para transportarla a una profundidad de \(10.3\,\text{m}\), permaneciendo constantes su temperatura, volumen y presión. Este requerimiento de energía es mayor que la energía que puede obtenerse mediante el metabolismo de alimentos con el oxígeno en esa cantidad de aire.
    \vspace*{8cm}
\end{itemize}

\end{document}
